\subsection{An integral rotation obstruction at multiplicity five}
\label{sec:1204-exclusion}

The ring $\RR{5}{1204}$ has no nontrivial proper fusion subring. Its
simple dimensions are
\[
 1,\quad3+2\sqrt3,\quad3+2\sqrt3,\quad7+4\sqrt3,
 \quad6+4\sqrt3,\quad8+4\sqrt3.
\]
We combine a forced central product with integrality of cyclic-rotation
eigenspaces.

\begin{proposition}\label{prop:1204-finite-profiles}
Every complex categorification of $\RR{5}{1204}$ is spherical and has a
pseudo-unitary Galois conjugate. Its central induction matrix belongs,
up to permutation of central simples, to an explicitly certified list of
$171$ restriction profiles. Their central ranks lie between $22$ and $51$.

\end{proposition}

\begin{proof}
The unrestricted signed-sphericalization test of Section~\ref{sec:signed}
enumerates $512$ associative signed quotient tables, partitioned into
$16$ sign-gauge orbits. One orbit consists of gauges of the original
fusion tensor. The exact dimension-character tests reject all $46$
character components on the other orbit representatives. As in the
sphericality argument of Proposition~\ref{prop:frontier-galois-pu}, a
categorification must therefore admit a spherical structure. Among the
character components of the original ring, the only one without a zero
simple dimension is the Frobenius--Perron orbit. Conjugating to its
positive embedding gives pseudo-unitarity. The positive global dimension
is $336+192\sqrt3$.

We may consequently apply the central-induction criteria with positive
dimensions. Exhaustive low-column enumeration gives $39$ frames. Only
two admit the required character assignments. Both assignments on the
first frame are excluded by exact full-character twist equations. The
second frame has three identical low restriction columns; its six
assignments are precisely the permutations of those columns. Relabelling
the corresponding low simples identifies these six necessary-data
problems, so one representative suffices.

For this representative, write the restriction Gram equations, joint
Galois multiplicity equations, and low cubic-word equations as an
integer nonnegative system $An=b$. Here $n$ includes both high-column
counts and low fusion multiplicities. The supplement records $165$
exact bound certificates. Each upper bound $n_j\le B$ is proved by a
separating functional for the augmented system
\[
 An=b,\qquad n_j-s=B+1,\qquad n,s\ge0;
\]
its infeasibility and integrality exclude $n_j\ge B+1$. Replacing $n_j$
by $-n_j$ gives lower bounds in the same way. Every bound is checked
against the preceding certified system. Bounded affine-lattice
enumeration produces exactly $39800$
high restriction profiles. Requiring integral nonnegative low-word
multiplicities through degree six and nonnegative restriction remainders
leaves $1833$ profiles.

For all unordered pairs of nonunit low simples, apply
Theorem~\ref{thm:joint-restriction-word-gram}. The profiles determine
$43$ distinct joint restriction--word Gram matrices. The exact fixed-column
Schur bounds of Lemma~\ref{lem:exact-rank-one-gram} enumerate every possible
joint coefficient column. The resulting count, Gram, and restriction
systems exclude $1568$ profiles by exact separating functionals. The
height-three full-character twist systems exclude another $91$.
Two additional joint count systems have inconsistent parity equations.
Finally, adjoining all cubic words in the last three low simples gives
an exact joint Gram obstruction for one further profile. Thus precisely
\[
 1833-1568-91-2-1=171
\]
profiles remain in the certified necessary domain. None is discarded
because a numerical optimizer reports infeasibility or exhausts its
running-time budget.

The supplement reconstructs this complete finite domain and checks every
separating and congruence certificate in exact arithmetic.
\end{proof}

\begin{theorem}\label{thm:1204-excluded}
The ring $\RR{5}{1204}$ has no complex fusion categorification.
\end{theorem}
\begin{proof}
By Proposition~\ref{prop:1204-finite-profiles}, it suffices to exclude the
$171$ profiles in the pseudo-unitary setting. Write the simple basis as
$(1,X,Y,Z,W,V)$ in the displayed dimension order. Every profile has a
unique central simple $U$ with $F(U)=X+Y$. All base simples are self-dual,
so uniqueness gives $U\simeq U^*$. The tensor gives
\[
 F(U^2)=2\cdot1+2X+2Y+2Z+4W+4V.
\]
The central unit occurs exactly once. The low restriction columns are
\[
 (1,0,0,0,0,0),\quad(1,2,2,4,4,4),\quad(1,2,2,5,4,5),
\]
and three copies of $(1,1,1,1,2,2)$. After removing the central unit,
the remaining unit coordinate is one, so precisely one further low
simple occurs. Coordinatewise bounds force its column to be
$(1,1,1,1,2,2)$. The remaining high restriction is
\[
 X+Y+Z+2W+2V.
\]
Its value under $c\mapsto c_W-c_Z$ is one. In every retained profile,
all active high columns except the pure $W$ column have nonpositive
value under this functional. Thus a central simple restricting to $W$
is necessary. Exactly $72$ profiles lack this column, leaving $99$.

For each remaining profile, apply
Proposition~\ref{prop:round3-full-character} through order six in
$K=\QQ(\zeta_{24})$. The global dimension $336+192\sqrt3$ has norm
$2304$, so the Cauchy prime support is $\{2,3\}$.
Proposition~\ref{prop:ambient-trace-domain} gives the complete finite
image of joint normalized traces, represented at height six by roots
of orders dividing $288$ together with the zero tuple. This bounds
trace types, not actual twist orders. On the columns singled out by
Proposition~\ref{prop:saturated-conductor-support}, the actual twist
order divides $12$; only these columns are restricted to twelfth-root
tuples. All five nonunit second indicators are $+1$, as forced by
odd self-dual multiplicities in the tensor squares.

The necessary equations use the induction Gram identities, indicator
formulas, every divisor moment of cyclic rotations of orders three
through six, low-pair modular equations and joint Galois balance.
For a rotation $E$, group its eigenvalues by the tuple of normalized
relative traces of all powers occurring in these equations. The
variable for a tuple is the sum of the corresponding eigenspace
dimensions, hence a nonnegative integer. Normalizing traces changes
coefficients, not multiplicities. Consequently the exact equations
$Ax=b$, $x\ge0$, retain integral rotation-count variables even if
high-simple counts are relaxed to real numbers.

The supplement supplies a short integral-bound certificate for each
profile. To prove $x_j\le B$ for an integral rotation count, adjoin the
opposite branch $x_j\ge B+1$, encoded by a nonnegative slack, and check
an integer ray $y$ with $y^TA\ge0$ and $y^Tb<0$. A lower bound is proved
by excluding $x_j\le B-1$. Subsequent steps may use preceding certified
bounds. The $99$ profiles require $112$ such bounds: $86$ use one and
$13$ use two. A final exact ray excludes each resulting system, giving
$211$ rays altogether. All rounded variables are rotation counts.
The checker reconstructs the equations and the full $72+99=171$
partition, then verifies every dot product with integers; numerical
optimization is disabled. Thus all necessary profiles are impossible.
\end{proof}
